\documentclass[../main.tex]{subfiles}
\begin{document}

\chapter{Introduction}
\lessoninfo{
  video={Lesson 1, videos 1--19},
  textbook={H\&P \cite{hennessy2017} §1.1--1.6},
  keywords={computer architecture, Moore's law, memory wall, dynamic power, static power, fabrication cost, yield},
}

This lesson defines what computer architecture is, explains why it is a
discipline in its own right, and introduces the technology trends and the
three quantities---performance, power and cost---that constrain every design
decision in the rest of the course.

\section{What is computer architecture?}

The discipline of designing a computer that is well suited to its purpose,
using the fabrication technology that is available, is called
\term{computer architecture}[コンピュータアーキテクチャ].  Two points in this
definition deserve emphasis.\source{Lesson 1, videos 2--3; H\&P §1.1--1.3.}

\begin{itemize}
  \item \emph{Well suited to its purpose.}  A mobile phone, a laptop and a
    data-center server are all computers, but they optimize for different
    things: battery life, size and weight in one case; raw performance and
    throughput per dollar in another.  Architecture is therefore about
    \emph{trade-offs}, not about a single best design.
  \item \emph{Available technology.}  The architect does not design
    transistors, but must know what the current process can deliver---how
    many transistors fit on a chip, how fast they switch, how much power they
    dissipate---because those numbers bound what is achievable.\margin{The
    scale involved: 2300 transistors on the Intel 4004 (1971), 80 billion on
    NVIDIA's H100 (2022).  Vendor data.}
\end{itemize}

\subsection{Why do we need it?}

Computer architecture serves two goals.\margin{The lecture summarizes these
as \enquote{improve performance} and \enquote{improve abilities}.}
\begin{enumerate}
  \item \textbf{Improve performance}, in the broad sense: execution speed,
    battery life, size, weight, energy efficiency.
  \item \textbf{Improve abilities}: make the computer able to do things it
    could not do before, e.g.\ 3D graphics, hardware support for debugging,
    security features.
\end{enumerate}

A further reason is timing.  A processor that is designed today will be
fabricated several years from now, so the architect must \emph{anticipate}
where the technology will be at that time rather than design for today's
parameters.\tip{Whenever a design decision in later lessons looks
over-engineered, ask which technology trend it was anticipating.}

\section{Technology trends}

\subsection{Moore's law}

\term{Moore's law}[ムーアの法則] is the empirical observation that the number
of transistors that can be placed on a chip of a given area doubles roughly
every 18 to 24 months.\source{\cite{moore1965}; Lesson 1, video 6; H\&P
§1.4.}  It is a trend, not a law of physics, but it held for several decades
and shaped the expectations of the whole industry.  The lecture derives the
following consequences, all with the same 18--24-month doubling
period:\caution{Moore said one year in 1965 and two years in 1975; the
popular \enquote{18 months} is David House's figure, not his.}
\begin{itemize}
  \item processor \emph{speed} doubles (more transistors allow more
    parallelism and smaller, faster gates);
  \item the \emph{energy per operation} halves;
  \item \emph{memory capacity} doubles.
\end{itemize}

\begin{example}
  If speed doubles every two years, a factor-of-$2^{5} = 32$ improvement is
  expected over a decade.  Conversely, a design team that spends four years
  producing a processor only $2\times$ faster than the one it replaces has
  merely kept pace with the trend.
\end{example}

\subsection{The memory wall}

Not every component follows the same curve.  While processor performance
doubled every 18--24 months, memory \emph{latency} improved only by about
$1.1\times$ over the same interval.\source{\cite{wulf1995}; Lesson 1,
video 8; H\&P §2.1.}  The result is a gap between how fast the processor
can consume data and how fast memory can supply it that widens every
generation.  This gap is called the \term{memory wall}[メモリウォール] and
is the reason why caches (Lessons 12--14) are so important.

\begin{figure}[htbp]
  \centering
  % Processor vs. memory performance trend (schematic, log scale).
\begin{tikzpicture}[x=1cm,y=1cm,>={Stealth[length=2mm]}, font=\small\sffamily]
  \draw[->] (0,0) -- (6.4,0) node[below left] {\iflangja{年}{year}};
  \draw[->] (0,0) -- (0,4.2) node[left] {\iflangja{性能（対数）}{performance (log)}};
  \draw[thick,ln-accent] (0.3,0.3) -- (6,3.9) node[pos=0.75,above left] {\iflangja{プロセッサ}{processor}};
  \draw[thick,ln-caution] (0.3,0.3) -- (6,1.2) node[pos=0.8,below] {\iflangja{メモリ（遅延）}{memory (latency)}};
  \draw[<->,ln-muted] (5.6,1.15) -- (5.6,3.65) node[midway,right,font=\footnotesize] {\iflangja{差}{gap}};
\end{tikzpicture}

  \caption{Processor performance and memory latency improve at different
    rates; the difference is the memory wall.}
  \label{fig:memory-wall}
\end{figure}

\begin{keypoint}
  Memory capacity keeps up with Moore's law; memory \emph{latency} does not.
  Most of the memory-system design in this course exists to hide that
  latency.\margin{On Skylake an L1 hit costs 4 cycles and an L2 hit 12;
  a DRAM access takes 60--\SI{90}{\nano\second}, i.e.\ 200--300 cycles at
  \SI{3}{\giga\hertz}.  Intel optimization manual; 7-cpu.com.}
\end{keypoint}

\section{Speed, power and cost}

For a long time \enquote{better} simply meant \enquote{faster}.  Today the
architect has to balance three quantities:\source{Lesson 1, videos 9--10.}
\begin{description}
  \item[Performance] how fast the machine runs its intended workload;
  \item[Power] how much energy it consumes, which sets battery life for
    portable devices and cooling cost for servers;\margin{Budgets: phone SoC
    2--\SI{5}{\watt}, desktop CPU 65--\SI{125}{\watt}, server CPU up to
    \SI{400}{\watt}.  Vendor data.}
  \item[Cost] how much it costs to fabricate.
\end{description}
Which of these dominates depends on the product: weight and battery life for
a phone, cost for an embedded controller, performance for a supercomputer.
\caution{\enquote{Faster} without a workload and a power/cost budget is not
a meaningful design goal.}

\section{Power consumption}

The power dissipated by a processor has two components.\source{Lesson 1,
videos 11--13; H\&P §1.5.}

\subsection{Dynamic (active) power}

The power consumed by charging and discharging the capacitance of wires and
gates whenever a signal switches is called \term{dynamic power}[動的電力]:
\begin{equation}
  P_{\mathrm{dyn}} = \tfrac{1}{2}\, C\, V^{2}\, f\, \alpha ,
  \label{eq:active-power}
\end{equation}
where $C$ is the total capacitance being switched, $V$ the supply voltage,
$f$ the clock frequency and $\alpha$ the \emph{activity factor}, the
fraction of the capacitance that actually switches each cycle.\margin{$C$
grows with chip size and transistor count; $\alpha$ depends on the workload
and on how much logic is left idle.  H\&P writes \enquote{frequency
switched} for the product $f\alpha$.}

Two facts follow directly from \cref{eq:active-power}.
\begin{itemize}
  \item Power scales with the \emph{square} of the voltage, so lowering $V$
    is the most effective lever.  However, the maximum frequency a circuit
    can sustain also drops as $V$ drops, so in practice $V$ and $f$ are
    lowered together and power falls roughly with $V^{3}$.
  \item Doubling the frequency at constant voltage doubles the power; the
    same speedup obtained by doubling the transistor count (more parallel
    hardware) costs power only through the larger $C$.\sidenote{Until about
    2005 each new process shrank $V$ together with the transistor, so a chip
    ran faster at constant power.  Leakage stopped voltage scaling near
    \SI{1}{\volt}: clock rates stalled at 3--\SI{4}{\giga\hertz} and the extra
    transistors went into more cores (Lesson 18).}
\end{itemize}

\begin{example}
  Reducing the supply voltage by \SI{15}{\percent} while keeping $C$, $f$ and
  $\alpha$ fixed reduces dynamic power to $0.85^{2} \approx 0.72$ of its
  original value, i.e.\ by about \SI{28}{\percent}.
\end{example}

\subsection{Static power}

The power consumed even when no signal switches is called
\term{static power}[静的電力] (leakage power); it arises because transistors are imperfect switches and leak current
when off.  Leakage grows as the supply voltage---and with it the transistor
threshold voltage---is reduced.  There is therefore a trade-off: lowering
$V$ reduces dynamic power but increases static power, and the total
$P = P_{\mathrm{dyn}} + P_{\mathrm{static}}$ has a minimum at some
intermediate voltage (\cref{fig:power}).

\begin{marginfigure}
  \resizebox{\linewidth}{!}{% Active vs. static power as a function of supply voltage (schematic).
\begin{tikzpicture}[x=1cm,y=1cm,>={Stealth[length=2mm]}, font=\small\sffamily]
  \draw[->] (0,0) -- (6.2,0) node[below left] {\iflangja{電源電圧 $V$}{supply voltage $V$}};
  \draw[->] (0,0) -- (0,4.2) node[left] {\iflangja{電力}{power}};
  % active ~ V^2
  \draw[thick,ln-accent] plot[domain=0.6:5.6,samples=40] (\x,{0.11*\x*\x});
  \node[ln-accent,anchor=west] at (4.3,3.4) {\iflangja{動的電力 $\propto V^2$}{active $\propto V^2$}};
  % static: grows as V drops (leakage at low threshold)
  \draw[thick,ln-caution] plot[domain=0.6:5.6,samples=40] (\x,{1.6/\x});
  \node[ln-caution,anchor=west] at (0.9,2.9) {\iflangja{静的電力}{static (leakage)}};
  % total
  \draw[thick,dashed,ln-muted] plot[domain=0.6:5.6,samples=40] (\x,{0.11*\x*\x+1.6/\x});
  \node[ln-muted,anchor=south west] at (4.0,1.1) {\iflangja{合計}{total}};
  \draw[dotted] (1.94,0) -- (1.94,{0.11*1.94*1.94+1.6/1.94});
  \node[below,font=\footnotesize] at (1.94,0) {$V_{\mathrm{opt}}$};
\end{tikzpicture}
}
  \caption{Dynamic and static power pull in opposite directions as the
    supply voltage changes.}
  \label{fig:power}
\end{marginfigure}

\begin{keypoint}[Power]
  Dynamic power $\propto C V^{2} f \alpha$ favors low voltage; static power
  penalizes it.  Modern designs sit near the minimum of the sum and get
  further gains from reducing $C$, $\alpha$ and $f$ rather than $V$.
\end{keypoint}

\section{Fabrication cost}

Chips are manufactured on circular silicon \term{wafer}[ウェハ]s, each of
which is cut into many identical rectangular \term{die}[ダイ]s.  The cost of
processing a wafer is roughly fixed, so the cost of a single die
is\source{Lesson 1, videos 15--17; H\&P §1.6.}
\begin{equation}
  \text{cost per die} \;=\;
  \frac{\text{cost per wafer}}
       {\text{dies per wafer} \times \text{yield}} .
  \label{eq:die-cost}
\end{equation}

\subsection{Yield}

Not every die works: defects are scattered over the wafer, and a die that
contains a defect is usually discarded.  The \term{yield}[歩留まり] is the
fraction of dies that are functional.  Because defects fall at random
positions, a larger die is more likely to contain one, so yield
\emph{decreases} as die area increases.\margin{H\&P §1.6 gives an empirical
model, $\text{yield} \approx 1/(1 + \text{defect density}\times
\text{area})^{N}$; a mature process leaves about \num{0.1} defects per
\si{\square\centi\metre}.}

\subsection{Consequences}

Combining \cref{eq:die-cost} with the yield behaviour gives the central
result of this section: halving the die area \emph{more than} halves the
cost per die, because it doubles the number of dies per wafer \emph{and}
raises the yield.\tip{Dies per wafer $\approx \pi (d/2)^{2}/A - \pi
d/\sqrt{2A}$; the second term removes the partial dies at the wafer edge.
H\&P §1.6.}  Read together with Moore's law this offers the architect
two options each technology generation:
\begin{itemize}
  \item keep the die area constant and pack in twice the transistors---a
    more capable chip at roughly the same cost;\margin{A \SI{300}{\milli\metre}
    wafer offers \SI{707}{\square\centi\metre}, and the largest die one
    lithography field can print is $26 \times \SI{33}{\milli\metre} =
    \SI{858}{\square\milli\metre}$; NVIDIA's H100 uses
    \SI{814}{\square\milli\metre} of it.} or
  \item keep the design constant and let it shrink to half the area---the
    same chip at well under half the cost.
\end{itemize}

\begin{example}
  Suppose a wafer costs \$5000 and holds 100 dies of a given design with a
  yield of \SI{50}{\percent}: the cost per working die is
  $5000/(100\times 0.5) = \$100$.  A shrink that fits 200 dies per wafer and
  raises the yield to \SI{70}{\percent} lowers the cost to
  $5000/(200\times 0.7) \approx \$36$, less than half.
\end{example}

\begin{keypoint}[Lesson summary]
  Architecture is trade-off engineering under three constraints---
  performance, power and cost---driven by technology trends (Moore's law) and
  their exceptions (the memory wall).  The next lesson makes
  \enquote{performance} precise.
\end{keypoint}

\end{document}
